\documentclass[11pt]{article}
\usepackage[utf8]{inputenc}
\usepackage{amsmath, amssymb, amsfonts}
\usepackage{geometry}
\geometry{a4paper, margin=1in}
\usepackage{xcolor}
\usepackage{listings}
\usepackage{hyperref}

% Python code formatting setup
\definecolor{codegreen}{rgb}{0,0.6,0}
\definecolor{codegray}{rgb}{0.5,0.5,0.5}
\definecolor{codepurple}{rgb}{0.58,0,0.82}
\definecolor{backcolour}{rgb}{0.95,0.95,0.92}

\lstdefinestyle{mystyle}{
    backgroundcolor=\color{backcolour},   
    commentstyle=\color{codegreen},
    keywordstyle=\color{magenta},
    numberstyle=\tiny\color{codegray},
    stringstyle=\color{codepurple},
    basicstyle=\ttfamily\footnotesize,
    breakatwhitespace=false,         
    breaklines=true,                 
    captionpos=b,                    
    keepspaces=true,                 
    numbers=none, % NO LINE NUMBERS AS REQUESTED
    showspaces=false,                
    showstringspaces=false,
    showtabs=false,                  
    tabsize=4
}
\lstset{style=mystyle}

\title{\textbf{Ultimate Preparation Guide: Continuous Fourier Series}}
\author{Part 1: Properties, Code, and Intuition}
\date{}

\begin{document}

\maketitle

\noindent \textbf{Prerequisite Utilities:} Before we begin, you will need a standard way to compute the Mean Squared Error (MSE) and extract coefficients into a predictable NumPy array format for easy math. Add these to your script tomorrow:

\begin{lstlisting}[language=Python]
import numpy as np

def calculate_mse(arr1, arr2):
    """Computes the Mean Squared Error between two complex arrays."""
    return np.mean(np.abs(arr1 - arr2)**2)

def get_coeff_array(fs_obj):
    """Extracts the dictionary of coefficients into a sorted 1D complex array."""
    # Sort keys to ensure we always go from -N to N in order
    sorted_keys = sorted(fs_obj.coeffs.keys())
    return np.array([fs_obj.coeffs[k] for k in sorted_keys], dtype=complex)
\end{lstlisting}

\hrulefill

\section{Layer 1: Basic Implementations (Single Properties)}

\subsection{Question 1: Linearity and DC Shift}
\textbf{Problem:} Given a signal $z(t)$, create a new signal $y(t) = A \cdot z(t) + (x_0 + j y_0)$. Verify the Linearity property of the Fourier Series.

\textbf{Theory \& Intuition:} 
Linearity states that if $x(t) \leftrightarrow a_k$ and $y(t) \leftrightarrow b_k$, then $A x(t) + B y(t) \leftrightarrow A a_k + B b_k$. 
Adding a constant $(x_0 + j y_0)$ is equivalent to adding a DC signal (frequency = 0). Therefore, \textit{only} the $c_0$ coefficient should change by exactly that constant. All other coefficients are simply scaled by $A$.

\begin{lstlisting}[language=Python]
# --- Setup ---
A = 2.5
dc_shift = 3.0 + 4.0j
y_signal = A * z + dc_shift

fs_y = FourierEpicycles(t, y_signal, N_HARMONICS)
fs_y.calculate_all_coefficients()

# --- Verification ---
c_y = get_coeff_array(fs_y)
c_z = get_coeff_array(fs) # Assuming fs is the original object

# Theoretical prediction
c_theoretical = A * c_z
# Find the index of the 0th harmonic (middle of the array)
zero_idx = N_HARMONICS 
c_theoretical[zero_idx] += dc_shift

mse_val = calculate_mse(c_y, c_theoretical)
print(f"Q1 Linearity MSE: {mse_val}")
\end{lstlisting}
\textbf{Code Explanation:} We scale the original signal and add a complex constant. In the theoretical array, we scale all original coefficients by $A$, and manually add the `dc_shift` \textit{only} to the $n=0$ index.

\subsection{Question 2: Time Shifting}
\textbf{Problem:} Shift the signal $z(t)$ in time by $t_0 = T/4$. Verify the Time Shifting property.

\textbf{Theory \& Intuition:} 
$x(t - t_0) \leftrightarrow c_n e^{-jn\omega t_0}$. 
Intuition: If you delay a signal, you aren't changing its shape, so the \textit{magnitudes} of the frequencies ($|c_n|$) stay exactly the same. However, to reconstruct the delayed shape, every rotating epicycle must start at a different angle. Higher frequencies spin faster, so they need a larger phase shift to account for the same time delay $t_0$. Hence, the phase shift $-n\omega t_0$ depends on $n$.

\begin{lstlisting}[language=Python]
# --- Setup ---
t0 = fs.T / 4.0

# Method 1: Using np.interp (Safest for continuous periodic signals)
# We evaluate z at (t - t0). We use modulo fs.T to handle wraparound.
t_shifted = (t - t0) % fs.T
# np.interp requires strictly increasing x-values, so we must separate real/imag
y_real = np.interp(t_shifted, t, z.real)
y_imag = np.interp(t_shifted, t, z.imag)
y_signal = y_real + 1j * y_imag

# Alternative Method 2: Using np.roll (Only works if t0 is an exact multiple of dt)
# shift_idx = int(t0 / (t[1] - t[0]))
# y_signal = np.roll(z[:-1], shift_idx)
# y_signal = np.append(y_signal, y_signal[0]) # restore closed interval

fs_y = FourierEpicycles(t, y_signal, N_HARMONICS)
fs_y.calculate_all_coefficients()

# --- Verification ---
c_y = get_coeff_array(fs_y)
c_z = get_coeff_array(fs)

n_arr = np.arange(-N_HARMONICS, N_HARMONICS + 1)
c_theoretical = c_z * np.exp(-1j * n_arr * fs.omega * t0)

print(f"Q2 Time Shift MSE: {calculate_mse(c_y, c_theoretical)}")
\end{lstlisting}
\textbf{Code Explanation:} `np.interp` is the most robust way to shift a signal in time because it doesn't rely on $t_0$ perfectly aligning with your array indices. We use modulo `% fs.T` to ensure the time values wrap around properly (since the signal is periodic).

\subsection{Question 3: Time Reversal and Conjugation}
\textbf{Problem:} Create a time-reversed signal $y(t) = z(-t)$ and a conjugated signal $w(t) = z^*(t)$. Verify their properties.

\textbf{Theory \& Intuition:} 
Time Reversal: $x(-t) \leftrightarrow c_{-n}$. Playing a signal backwards reverses the direction of all the rotating epicycles.
Conjugation: $x^*(t) \leftrightarrow c_{-n}^*$. Conjugating the spatial coordinates flips the shape over the x-axis.

\begin{lstlisting}[language=Python]
# --- Setup ---
# Time Reversal: Reverse the array. 
# Because t[0]=0 and t[-1]=T are the same point, we reverse the inner elements.
y_signal = np.zeros_like(z)
y_signal[0] = z[0]
y_signal[1:] = z[:0:-1] # slice from end down to index 1

# Conjugation
w_signal = np.conj(z)

fs_y = FourierEpicycles(t, y_signal, N_HARMONICS)
fs_y.calculate_all_coefficients()
fs_w = FourierEpicycles(t, w_signal, N_HARMONICS)
fs_w.calculate_all_coefficients()

# --- Verification ---
c_y = get_coeff_array(fs_y)
c_w = get_coeff_array(fs_w)
c_z = get_coeff_array(fs)

# Theoretical: c_{-n} is just the reversed coefficient array!
c_theoretical_y = c_z[::-1] 
c_theoretical_w = np.conj(c_z[::-1])

print(f"Q3 Time Reversal MSE: {calculate_mse(c_y, c_theoretical_y)}")
print(f"Q3 Conjugation MSE: {calculate_mse(c_w, c_theoretical_w)}")
\end{lstlisting}

\subsection{Question 4: Frequency Shifting (Modulation)}
\textbf{Problem:} Multiply the signal by a complex exponential $y(t) = e^{j M \omega t} z(t)$ where $M = 5$. Verify the property.

\textbf{Theory \& Intuition:} 
$e^{j M \omega t} x(t) \leftrightarrow c_{n-M}$. 
Intuition: Multiplying by a complex exponential in time "spins" the entire drawing. In the frequency domain, this shifts all the coefficients to the right by $M$ slots.

\begin{lstlisting}[language=Python]
# --- Setup ---
M = 5
y_signal = z * np.exp(1j * M * fs.omega * t)

fs_y = FourierEpicycles(t, y_signal, N_HARMONICS)
fs_y.calculate_all_coefficients()

# --- Verification ---
c_y = get_coeff_array(fs_y)
c_z = get_coeff_array(fs)

# Theoretical: Shift the array to the right by M. 
# Note: Because we have a finite N, shifting drops the last M elements 
# and we don't know the new first M elements. We verify only the overlapping part.
c_y_valid = c_y[M:]           # The computed shifted coefficients
c_theo_valid = c_z[:-M]       # The original coefficients

print(f"Q4 Frequency Shift MSE: {calculate_mse(c_y_valid, c_theo_valid)}")
\end{lstlisting}

\hrulefill

\section{Layer 2: Harder (Calculus, Scaling, \& Combinations)}

\subsection{Question 5: Time Scaling}
\textbf{Problem:} Compress the signal in time by a factor of $\alpha = 2$, so $y(t) = z(2t)$. Verify the Time Scaling property.

\textbf{Theory \& Intuition:} 
$x(\alpha t) \leftrightarrow c_n$. 
Intuition: Playing a video at 2x speed doesn't change the frames of the video (the shape, $c_n$), it only changes how fast they play (the fundamental frequency becomes $\alpha \omega$, and the period becomes $T/\alpha$). 

\begin{lstlisting}[language=Python]
# --- Setup ---
alpha = 2.0
# The new period is T / alpha
T_new = fs.T / alpha
# We must create a new time array for the new period
t_new = np.linspace(0, T_new, len(t))

# The signal values themselves don't change, they just happen over a shorter time!
y_signal = z.copy() 

# Initialize with the NEW time array
fs_y = FourierEpicycles(t_new, y_signal, N_HARMONICS)
fs_y.calculate_all_coefficients()

# --- Verification ---
c_y = get_coeff_array(fs_y)
c_z = get_coeff_array(fs)

# Theoretical: The coefficients are EXACTLY the same.
print(f"Q5 Time Scaling MSE: {calculate_mse(c_y, c_z)}")
# Also verify the frequency changed
print(f"Original Omega: {fs.omega}, New Omega: {fs_y.omega} (Should be {alpha * fs.omega})")
\end{lstlisting}
\textbf{Code Explanation:} This is a trick question regarding OOP. You do not interpolate the signal array. You simply pass the exact same `z` array into `FourierEpicycles`, but you pass a \textit{compressed time array} `t_new`. The class will automatically calculate the new $T$ and $\omega$.

\subsection{Question 6: Differentiation Property}
\textbf{Problem:} Compute the derivative of the signal $y(t) = \frac{d}{dt}z(t)$. Verify the differentiation property.

\textbf{Theory \& Intuition:} 
$\frac{dx(t)}{dt} \leftrightarrow j n \omega c_n$. 
Intuition: The derivative measures the rate of change. High-frequency epicycles spin faster, so their rate of change is much larger (hence multiplying by $n\omega$). The $j$ represents a $90^\circ$ phase shift (since the derivative of sine is cosine).

\begin{lstlisting}[language=Python]
# --- Setup ---
# Numerical differentiation using np.gradient
# Since t is uniformly spaced, we can just pass the spacing dt
dt = t[1] - t[0]
y_signal = np.gradient(z, dt)

fs_y = FourierEpicycles(t, y_signal, N_HARMONICS)
fs_y.calculate_all_coefficients()

# --- Verification ---
c_y = get_coeff_array(fs_y)
c_z = get_coeff_array(fs)

n_arr = np.arange(-N_HARMONICS, N_HARMONICS + 1)
c_theoretical = 1j * n_arr * fs.omega * c_z

print(f"Q6 Differentiation MSE: {calculate_mse(c_y, c_theoretical)}")
\end{lstlisting}
\textbf{Alternative Method:} Instead of `np.gradient`, you could use finite differences: `y_signal = np.diff(z) / dt`, but `np.diff` reduces the array size by 1. You would have to append a value to keep the shape consistent. `np.gradient` handles boundaries automatically and maintains the array shape, making it superior here.

\subsection{Question 7: Parseval's Relation (Energy Conservation)}
\textbf{Problem:} Verify Parseval's theorem for the signal $z(t)$.

\textbf{Theory \& Intuition:} 
$\frac{1}{T} \int_0^T |x(t)|^2 dt = \sum_{n=-\infty}^{\infty} |c_n|^2$. 
Intuition: The total energy of the signal in the time domain must exactly equal the sum of the energies of all its individual frequency components. (Note: Because we only compute up to $N$ harmonics, the frequency side will be slightly smaller than the time side. The MSE won't be exactly 0, but it should be very close if $N$ is large).

\begin{lstlisting}[language=Python]
# Time domain power (using np.trapezoid as required)
power_time = np.trapezoid(np.abs(z)**2, t) / fs.T

# Frequency domain power
c_z = get_coeff_array(fs)
power_freq = np.sum(np.abs(c_z)**2)

print(f"Q7 Parseval's - Time Power: {power_time:.4f}, Freq Power: {power_freq:.4f}")
print(f"Difference (Energy in truncated harmonics): {power_time - power_freq:.6f}")
\end{lstlisting}

\subsection{Question 8: Conjugate Symmetry and Even/Odd Decomposition}
\textbf{Problem:} Extract the real part of the signal $x(t) = \Re\{z(t)\}$. Verify that its coefficients exhibit conjugate symmetry ($c_n = c_{-n}^*$). Then, decompose $x(t)$ into Even and Odd parts and verify their coefficient properties.

\textbf{Theory \& Intuition:} 
If a signal is purely real, its negative frequencies must perfectly cancel out the imaginary parts of its positive frequencies. Therefore, $c_{-n}$ must be the complex conjugate of $c_n$. 
Furthermore, the Even part of a real signal corresponds to the Real part of $c_n$, and the Odd part corresponds to the Imaginary part of $c_n$.

\begin{lstlisting}[language=Python]
# --- Setup ---
x_real = z.real
fs_real = FourierEpicycles(t, x_real, N_HARMONICS)
fs_real.calculate_all_coefficients()
c_real = get_coeff_array(fs_real)

# 1. Verify Conjugate Symmetry: c_n == conj(c_{-n})
# c_real[::-1] gives us the array reversed (which maps n to -n)
mse_sym = calculate_mse(c_real, np.conj(c_real[::-1]))
print(f"Q8 Conjugate Symmetry MSE: {mse_sym}")

# 2. Even/Odd Decomposition
# x_even(t) = (x(t) + x(-t)) / 2
x_rev = np.zeros_like(x_real)
x_rev[0] = x_real[0]
x_rev[1:] = x_real[:0:-1]

x_even = (x_real + x_rev) / 2.0
x_odd = (x_real - x_rev) / 2.0

fs_even = FourierEpicycles(t, x_even, N_HARMONICS)
fs_even.calculate_all_coefficients()
fs_odd = FourierEpicycles(t, x_odd, N_HARMONICS)
fs_odd.calculate_all_coefficients()

c_even = get_coeff_array(fs_even)
c_odd = get_coeff_array(fs_odd)

# Theoretical: c_even should be purely real, c_odd purely imaginary
# We compare them to the real and imaginary parts of the original c_real
mse_even = calculate_mse(c_even, c_real.real)
mse_odd = calculate_mse(c_odd, 1j * c_real.imag)

print(f"Q8 Even Part MSE: {mse_even}")
print(f"Q8 Odd Part MSE: {mse_odd}")
\end{lstlisting}

\hrulefill

\section{Layer 3: Very Hard (Convolutions, Integrations, Multi-layered)}

\subsection{Question 9: Periodic Convolution}
\textbf{Problem:} Given two signals $x(t)$ and $y(t)$ (let's use $z(t)$ and a shifted version of $z(t)$), compute their periodic convolution $w(t) = \int_0^T x(\tau) y(t-\tau) d\tau$. Verify the convolution property.

\textbf{Theory \& Intuition:} 
$x(t) \circledast y(t) \leftrightarrow T \cdot c_n^{(x)} c_n^{(y)}$. 
Convolution in the time domain is equivalent to multiplication in the frequency domain (scaled by $T$). Because you cannot use `np.fft` or `scipy.signal.convolve` (which uses FFT under the hood), you must implement the convolution integral manually using `np.trapezoid`.

\begin{lstlisting}[language=Python]
# --- Setup ---
x_sig = z
# Let y_sig be a simple sine wave to make convolution fast and clean
y_sig = np.sin(fs.omega * t) 

# Manual O(N^2) Periodic Convolution using np.trapezoid
w_sig = np.zeros_like(t, dtype=complex)
# We need y(-tau). 
y_rev = np.zeros_like(y_sig)
y_rev[0] = y_sig[0]
y_rev[1:] = y_sig[:0:-1]

for i in range(len(t)):
    # Shift y(-tau) to the right by i steps to get y(t - tau)
    y_shifted = np.roll(y_rev, i)
    # Integrate x(tau) * y(t - tau) d_tau
    w_sig[i] = np.trapezoid(x_sig * y_shifted, t)

fs_y = FourierEpicycles(t, y_sig, N_HARMONICS)
fs_y.calculate_all_coefficients()
fs_w = FourierEpicycles(t, w_sig, N_HARMONICS)
fs_w.calculate_all_coefficients()

# --- Verification ---
c_x = get_coeff_array(fs)
c_y = get_coeff_array(fs_y)
c_w = get_coeff_array(fs_w)

c_theoretical = fs.T * c_x * c_y

print(f"Q9 Periodic Convolution MSE: {calculate_mse(c_w, c_theoretical)}")
\end{lstlisting}
\textbf{Code Explanation:} We manually implement the convolution integral. `np.roll(y_rev, i)` perfectly simulates $y(t - \tau)$ because the signal is periodic and uniformly sampled. We integrate the product using `np.trapezoid`.

\subsection{Question 10: Multiplication in Time}
\textbf{Problem:} Multiply two signals in the time domain $w(t) = x(t)y(t)$. Verify the property.

\textbf{Theory \& Intuition:} 
$x(t)y(t) \leftrightarrow \sum_{l=-\infty}^{\infty} c_l^{(x)} c_{n-l}^{(y)}$. 
Multiplication in time is convolution in frequency! The coefficients of the multiplied signal are the discrete convolution of the individual coefficient arrays.

\begin{lstlisting}[language=Python]
# --- Setup ---
x_sig = z
y_sig = np.cos(2 * fs.omega * t) # A 2nd harmonic cosine
w_sig = x_sig * y_sig

fs_y = FourierEpicycles(t, y_sig, N_HARMONICS)
fs_y.calculate_all_coefficients()
fs_w = FourierEpicycles(t, w_sig, N_HARMONICS)
fs_w.calculate_all_coefficients()

# --- Verification ---
c_x = get_coeff_array(fs)
c_y = get_coeff_array(fs_y)
c_w = get_coeff_array(fs_w)

# Discrete convolution of coefficients. 
# np.convolve does exactly sum(a[l] * b[n-l]). 
# Mode 'same' keeps the array size 2N+1.
c_theoretical = np.convolve(c_x, c_y, mode='same')

print(f"Q10 Multiplication MSE: {calculate_mse(c_w, c_theoretical)}")
\end{lstlisting}

\subsection{Question 11: The Integration Property}
\textbf{Problem:} Compute the integral of the signal $y(t) = \int_{-\infty}^t x(\tau) d\tau$. Verify the property.

\textbf{Theory \& Intuition:} 
$\int x(t) dt \leftrightarrow \frac{c_n}{jn\omega}$. 
Integration is the opposite of differentiation, so we divide by $jn\omega$. \textbf{CRITICAL CATCH:} This property is \textit{only} valid if $c_0 = 0$ (the signal has zero mean). If $c_0 \neq 0$, integrating a constant yields a ramp function ($c_0 t$), which goes to infinity and is no longer periodic!

\begin{lstlisting}[language=Python]
import scipy.integrate

# --- Setup ---
# 1. We MUST remove the DC component first!
c0 = fs.coeffs[0]
x_zero_mean = z - c0

# 2. Integrate using cumulative trapezoid
# initial=0 ensures the output array is the same size as t
y_signal = scipy.integrate.cumulative_trapezoid(x_zero_mean, t, initial=0)

fs_y = FourierEpicycles(t, y_signal, N_HARMONICS)
fs_y.calculate_all_coefficients()

# --- Verification ---
c_y = get_coeff_array(fs_y)
c_x = get_coeff_array(fs) # Original coefficients

n_arr = np.arange(-N_HARMONICS, N_HARMONICS + 1)
c_theoretical = np.zeros_like(c_x, dtype=complex)

# Apply formula ONLY for n != 0 to avoid division by zero
for i, n in enumerate(n_arr):
    if n != 0:
        c_theoretical[i] = c_x[i] / (1j * n * fs.omega)
    else:
        # The DC component of the integrated signal depends on the constant of integration.
        # We just copy the computed one to ignore it in the MSE calculation.
        c_theoretical[i] = c_y[i] 

print(f"Q11 Integration MSE: {calculate_mse(c_y, c_theoretical)}")
\end{lstlisting}
\textbf{Code Explanation:} We use `scipy.integrate.cumulative_trapezoid` to perform the running integral. We must manually handle the $n=0$ case in the theoretical array to avoid a `ZeroDivisionError`.

\subsection{Question 12: The Ultimate Combo}
\textbf{Problem:} Given $z(t)$, construct $y(t) = \frac{d}{dt} [z(t - t_0)] \cdot e^{j \omega_0 t}$. Find the theoretical coefficients and verify.

\textbf{Theory \& Intuition:} 
Apply properties sequentially:
1. Shift: $c_n^{(1)} = c_n e^{-jn\omega t_0}$
2. Differentiate: $c_n^{(2)} = jn\omega \cdot c_n^{(1)} = jn\omega c_n e^{-jn\omega t_0}$
3. Frequency Shift (by $M = \omega_0 / \omega$): $c_n^{(3)} = c_{n-M}^{(2)} = j(n-M)\omega c_{n-M} e^{-j(n-M)\omega t_0}$.

\begin{lstlisting}[language=Python]
# --- Setup ---
t0 = fs.T / 8.0
M = 2

# 1. Shift
t_shifted = (t - t0) % fs.T
z_shift = np.interp(t_shifted, t, z.real) + 1j * np.interp(t_shifted, t, z.imag)
# 2. Differentiate
dt = t[1] - t[0]
z_diff = np.gradient(z_shift, dt)
# 3. Frequency Shift
y_signal = z_diff * np.exp(1j * M * fs.omega * t)

fs_y = FourierEpicycles(t, y_signal, N_HARMONICS)
fs_y.calculate_all_coefficients()

# --- Verification ---
c_y = get_coeff_array(fs_y)
c_z = get_coeff_array(fs)
n_arr = np.arange(-N_HARMONICS, N_HARMONICS + 1)

# Theoretical calculation
c_theo_full = 1j * n_arr * fs.omega * c_z * np.exp(-1j * n_arr * fs.omega * t0)
# Shift the theoretical array by M
c_theo_valid = c_theo_full[:-M]
c_y_valid = c_y[M:]

print(f"Q12 Ultimate Combo MSE: {calculate_mse(c_y_valid, c_theo_valid)}")
\end{lstlisting}

\end{document}